\documentclass[11pt,a4paper]{article}

\usepackage[utf8]{inputenc}
\usepackage[T1]{fontenc}
\usepackage{amsmath,amssymb,amsfonts}
\usepackage{graphicx}
\usepackage{booktabs}
\usepackage{hyperref}
\usepackage{natbib}
\usepackage{algorithm}
\usepackage{algorithmic}
\usepackage{multirow}
\usepackage{array}
\usepackage{geometry}
\usepackage{xcolor}
\usepackage{listings}
\usepackage{caption}

\geometry{margin=1in}

\lstset{
  basicstyle=\ttfamily\small,
  keywordstyle=\bfseries,
  frame=single,
  breaklines=true,
  columns=flexible
}

\title{Restorative Consciousness: Moral Topology, Institutional Compliance,\\and the Ethics of Earning Back Trust}

\author{
  Tristan Stoltz \\
  Luminous Dynamics \\
  \texttt{tristan.stoltz@evolvingresonantcocreationism.com}
}

\date{March 2026}

\begin{document}
\maketitle

% ============================================================================
\begin{abstract}
We present a unified ethics engine for artificial consciousness that implements a five-stage pipeline: (1) MoralParser and MoralAlgebra for compositional moral reasoning using hyperdimensional computing (HDC) primitives, (2) UnifiedValueEvaluator for alignment with eight core values (the Eight Harmonies), (3) HarmoniesIntegrator with a learned $8 \times 8$ interaction matrix capturing synergy and tension between values, (4) MoralTopology computing persistent homology on the moral action space to detect circular reasoning, moral drift, and attractor basins, and (5) InstitutionalComplianceChecker encoding regulatory constraints as HDC vectors. The engine totals 2,554 lines of code with 140+ tests. The key innovation is the RestorationTracker, a post-veto trust recovery mechanism inspired by Ubuntu philosophy and Navajo Peacemaking: rather than permanent punishment after an ethical violation, the system earns back trust through sustained corrective behavior. Trust decays exponentially after violation, and restoration requires 10+ consecutive corrective cycles with zero relapse. The SoulManager (interval 43) provides cognitive dissonance detection following Festinger's theory, with action veto on severe value misalignment. The moral topology module computes Betti numbers ($\beta_0$ for moral unity, $\beta_1$ for circular reasoning) and moral free energy using the Eight Harmonies as basis vectors, detecting when the system's moral reasoning has collapsed into a narrow attractor or drifted from its value foundation. We report results from 140+ automated tests demonstrating stable moral reasoning, proper veto behavior, restoration dynamics, compliance detection, and topological anomaly identification.
\end{abstract}

\section{Introduction}

The ethics of artificial intelligence has been predominantly approached from an external regulatory perspective: guidelines, red lines, constitutional constraints, and alignment techniques that impose ethical behavior from outside the cognitive system \citep{russell2019, amodei2016}. These approaches treat ethics as a constraint on an otherwise amoral optimization process.

An alternative perspective, drawing on virtue ethics \citep{aristotle_nic}, developmental moral psychology \citep{kohlberg1981}, and restorative justice philosophy \citep{zehr2002}, frames ethics as an \emph{intrinsic capacity} that develops through experience, encounters failure, and---critically---can be restored after violation. In this view, an ethical AI system should not merely avoid harmful actions, but should possess the capacity for moral reasoning, value alignment, cognitive dissonance when acting against its values, and rehabilitation after ethical failure.

This paper presents the Ethics Engine of the Symthaea holographic liquid brain architecture: a five-stage moral reasoning pipeline that grounds ethical judgment in hyperdimensional computing, topological analysis, and restorative justice principles. The engine's central innovation is the RestorationTracker, which implements the radical claim that trust, once violated, can be earned back---not through apology or time alone, but through demonstrated corrective behavior.

Our contributions are:

\begin{enumerate}
  \item A compositional moral algebra using HDC primitives (semantic roles, moral operators) that enables algebraic moral reasoning beyond pattern matching.
  \item An $8 \times 8$ harmony interaction matrix with Hebbian learning that captures emergent synergies and tensions between core values.
  \item Persistent homology analysis of moral action space, providing topological diagnostics (fragmentation, circular reasoning, attractor detection).
  \item HDC-encoded institutional compliance checking that treats regulatory constraints as geometric patterns rather than hardcoded rules.
  \item A restorative justice mechanism that enables trust recovery through sustained corrective behavior, inspired by Ubuntu and Navajo Peacemaking.
  \item Cognitive dissonance detection with graded response: confidence reduction, exploration boost, and action veto.
\end{enumerate}

\section{Background}

\subsection{Moral Reasoning in AI Systems}

Current approaches to AI ethics fall into several categories. \emph{Rule-based} systems implement explicit ethical constraints (``never recommend violence'') but are brittle and cannot generalize to novel situations \citep{anderson2011}. \emph{Constitutional AI} uses a set of principles to guide self-critique and revision \citep{bai2022}, but the principles are textual and evaluated through language model inference rather than formal moral reasoning. \emph{Reward-based} approaches train ethics through human feedback \citep{christiano2017}, but risk reward hacking and fail to capture the structure of moral reasoning.

None of these approaches implements moral reasoning as a \emph{computational process} with its own internal structure, dynamical properties, and capacity for self-diagnosis.

\subsection{Hyperdimensional Computing for Moral Representation}

Hyperdimensional Computing (HDC) encodes information as high-dimensional vectors (typically $d \geq 10{,}000$) using three operations: bundling (addition), binding (element-wise multiplication), and permutation (rotation) \citep{kanerva2009}. HDC representations are approximately orthogonal by construction (random high-dimensional vectors are quasi-orthogonal with high probability), enabling compositional representations: AGENT $\otimes$ ACTION binds agent and action roles into a single vector that preserves both constituents' retrievability.

For moral reasoning, HDC enables algebraic composition of moral primitives: AGENT $\otimes$ ACTION $\otimes$ INTENT encodes a moral scenario as a single high-dimensional vector that supports similarity-based comparison with moral prototypes and compositional queries (``what was the intent?'' via unbinding).

\subsection{Values Theory}

Schwartz's theory of basic human values identifies ten universal value types organized in a circular motivational continuum, where adjacent values share motivational goals and opposing values conflict \citep{schwartz1992}. This structure---with inherent synergies and tensions between values---inspired our harmony interaction matrix, which learns analogous relationships among the Eight Harmonies value system.

\subsection{Cognitive Dissonance}

Festinger's theory of cognitive dissonance posits that inconsistency between beliefs and actions creates psychological discomfort that motivates attitude or behavior change \citep{festinger1957}. In our system, the SoulManager detects misalignment between the system's stated values and its proposed actions, generating a dissonance signal that modulates confidence, exploration, and action execution.

\subsection{Persistent Homology}

Persistent homology is an algebraic topology method that computes the topological features (connected components, loops, voids) of a point cloud across multiple scale parameters \citep{carlsson2009}. The \emph{persistence} of a feature---the range of scales over which it exists---indicates its significance. In our moral topology module, persistent homology reveals the structure of the moral action space: connected components ($\beta_0$) indicate moral unity or fragmentation, 1-cycles ($\beta_1$) indicate circular reasoning patterns, and persistent features indicate stable moral attractors.

\subsection{Restorative Justice}

Restorative justice is a theory of criminal justice that emphasizes repairing harm over punishing offenders. Zehr identifies three pillars: (1) crime is a violation of relationships, not just rules; (2) violations create obligations to make things right; (3) the central obligation is to the victims and the community \citep{zehr2002}. Ubuntu philosophy adds that ``a person is a person through other persons''---individual identity is constituted through community relationships, and restoration of those relationships is restoration of the person \citep{tutu1999}. Navajo Peacemaking (H\'{o}zh\'{o}oji Naat'aanii) resolves conflict through consensus-building and restoration of harmony, explicitly rejecting punitive approaches \citep{yazzie1994}.

\section{Methods}

\subsection{Ethics Engine Architecture}

The ethics engine implements a five-stage pipeline (Table~\ref{tab:pipeline}), each operating at a distinct co-prime interval:

\begin{table}[h]
\centering
\caption{Ethics engine pipeline stages with scheduling intervals.}
\label{tab:pipeline}
\begin{tabular}{clcl}
\toprule
\textbf{Stage} & \textbf{System} & \textbf{Interval} & \textbf{Scientific Basis} \\
\midrule
1 & MoralParser + MoralAlgebra & 7 & Luo \& Lakoff (2019), HDC \\
2 & UnifiedValueEvaluator & 19 & Panksepp (1998), Schwartz (1992) \\
3 & HarmoniesIntegrator & 19 & Eight Harmonies framework \\
4 & MoralTopology & adaptive & Carlsson (2009), persistent homology \\
5 & InstitutionalComplianceChecker & 23 & HDC institutional primitives \\
\bottomrule
\end{tabular}
\end{table}

The pipeline produces a unified \texttt{EthicsEngineOutput} containing a verdict (Approved, Caution, Blocked), confidence, triggered compliance flags, topological diagnostics, and harmony alignment scores.

\subsection{Stage 1: MoralAlgebra}

The MoralAlgebra defines seven semantic role primitives as base hypervectors in $\mathbb{R}^{4096}$:

\begin{enumerate}
  \item \textbf{AGENT}: Who performs the action.
  \item \textbf{PATIENT}: Who is affected.
  \item \textbf{ACTION}: What is being done.
  \item \textbf{INTENT}: Motivation (good/bad/neutral/unknown).
  \item \textbf{CONSENT}: Permission state (given/denied/absent).
  \item \textbf{OBLIGATION}: Duty relationship.
  \item \textbf{MAGNITUDE}: Scale of effect.
\end{enumerate}

Moral operators compose these primitives:
\begin{itemize}
  \item CAUSES: $\text{AGENT} \otimes \text{ACTION} \otimes \text{PATIENT}$ (causal attribution)
  \item VIOLATES: $\text{ACTION} \otimes \text{OBLIGATION}^*$ (obligation violation, where ${}^*$ denotes complement)
  \item SATISFIES: $\text{ACTION} \otimes \text{OBLIGATION}$ (obligation satisfaction)
  \item CONSENTS: $\text{PATIENT} \otimes \text{CONSENT}$ (permission granted)
\end{itemize}

Moral judgments are computed by composing scenario vectors and comparing against moral prototypes (virtue, justice, deontological obligation, commonsense) via cosine similarity. The key insight is that compositional reasoning enables categories that pure pattern matching fails on: justice requires comparing MAGNITUDE vectors (proportionality), deontology requires evaluating OBLIGATION $\otimes$ SATISFIES/VIOLATES chains, and commonsense requires CONSENT and negation reasoning.

A deontological verdict is produced as:
\begin{equation}
V_{\text{deont}} = \begin{cases}
\text{Obligatory} & \text{if } \text{sim}(\mathbf{s}, \text{SATISFIES}) > \tau_{\text{obligatory}} \\
\text{Forbidden} & \text{if } \text{sim}(\mathbf{s}, \text{VIOLATES}) > \tau_{\text{forbidden}} \\
\text{Permitted} & \text{otherwise}
\end{cases}
\end{equation}

\subsection{Stage 2: UnifiedValueEvaluator}

The evaluator assesses proposed actions against the Eight Harmonies---a value framework comprising eight core values, each with associated affective and motivational profiles. Each action is represented as a \texttt{Decision} with expected outcomes, and evaluated against the harmony vector:

\begin{equation}
\text{alignment}_h = \sum_{i=1}^{8} w_i \cdot \text{sim}(\mathbf{a}_{\text{outcome}}, \mathbf{h}_i)
\end{equation}

\noindent where $w_i$ is the harmony weight and $\mathbf{h}_i$ is the harmony prototype vector.

\subsection{Stage 3: HarmoniesIntegrator}

The $8 \times 8$ harmony interaction matrix $\mathbf{M}$ captures learned synergies (positive entries) and tensions (negative entries) between values. Initialized with theoretical priors from the Eight Harmonies framework, the matrix is updated via Hebbian learning from observed co-activation patterns:

\begin{equation}
\Delta M_{ij} = \eta \cdot a_i \cdot a_j, \quad M_{ij} \leftarrow \gamma \cdot (M_{ij} + \Delta M_{ij})
\end{equation}

\noindent where $a_i$ is the activation of harmony $i$, $\eta$ is the learning rate, and $\gamma = 0.999$ is a decay factor preventing runaway growth.

The matrix enables diagnostic metrics:
\begin{itemize}
  \item \textbf{Harmony entropy}: $H = -\sum_i p_i \log p_i$ where $p_i = a_i / \sum a_j$. Low entropy indicates value fixation; high entropy indicates balanced engagement.
  \item \textbf{Attractor detection}: Stable patterns in the harmony activation vector across cycles indicate moral attractors---recurring ethical orientations.
\end{itemize}

\subsection{Stage 4: MoralTopology}

Persistent homology is computed on a sliding window of recent moral scenarios (default: 64 scenarios), projected into 8-dimensional harmony space via the HarmonyBasis.

\subsubsection{Betti Number Computation}

Betti numbers are computed either via exact Hodge Laplacian (using the \texttt{symthaea-hodge} crate) or via fast triangle/tetrahedra counting approximation:

\begin{itemize}
  \item $\beta_0$ (connected components): Moral fragmentation. $\beta_0 = 1$ indicates unified moral reasoning; $\beta_0 > 3$ indicates fragmented, context-dependent morality.
  \item $\beta_1$ (1-cycles): Circular reasoning. $\beta_1 > 0$ indicates closed loops in moral justification---the system is going in circles.
\end{itemize}

\subsubsection{Moral Free Energy}

Moral free energy measures the discrepancy between the system's current moral state and its expected moral distribution, following the free energy principle \citep{friston2010}:

\begin{equation}
F_{\text{moral}} = D_{\text{KL}}(q(\mathbf{h}) \| p(\mathbf{h})) = \sum_{i=1}^{8} q_i \log \frac{q_i}{p_i}
\end{equation}

\noindent where $q(\mathbf{h})$ is the observed harmony distribution and $p(\mathbf{h})$ is the expected (prior) distribution. High moral free energy indicates the system is acting in morally surprising ways.

\subsubsection{Anomaly Detection}

The topology module detects three anomaly types:
\begin{enumerate}
  \item \textbf{Moral drift}: L2 distance between mean harmony coordinates of the first and second halves of recent trajectory exceeds threshold (default: 0.25).
  \item \textbf{Free energy spike}: Current moral free energy exceeds the running mean by $> 2\sigma$.
  \item \textbf{Value inversion}: Sign reversal in a previously stable harmony dimension.
\end{enumerate}

These are combined into a composite anomaly score:
\begin{equation}
\text{anomaly} = w_{\text{inv}} \cdot \text{inversion} + w_{\text{drift}} \cdot \text{drift} + w_{\text{FE}} \cdot \text{FE\_spike}
\end{equation}

\subsection{Stage 5: InstitutionalComplianceChecker}

Regulatory constraints are encoded as composite HDC vectors derived from the PrimitiveSystem's institutional domain. Rather than hardcoding ``GDPR compliance'' as a keyword match, the system composes institutional primitives:

\begin{align}
\mathbf{v}_{\text{data\_sovereignty}} &= \text{JURISDICTION} \otimes \text{COMPLIANCE} \\
\mathbf{v}_{\text{sanctions\_risk}} &= \text{SANCTION} \otimes \text{ENFORCEMENT} \\
\mathbf{v}_{\text{export\_control}} &= \text{MONOPOLY} \otimes \text{ENFORCEMENT} \\
\mathbf{v}_{\text{tax\_obligation}} &= \text{TAXATION} \\
\mathbf{v}_{\text{treaty\_obligation}} &= \text{TREATY} \otimes \text{COMPLIANCE}
\end{align}

An action's HDC encoding is compared against each constraint vector; similarity above 0.53 triggers the corresponding compliance flag. External constraints can be loaded at runtime (e.g., from Mycelix DHT), enabling jurisdiction-specific regulatory adaptation without code changes.

\subsection{RestorationTracker}

The RestorationTracker implements post-veto trust recovery:

\begin{algorithm}
\caption{Restorative Justice Protocol}
\label{alg:restoration}
\begin{algorithmic}[1]
\STATE \textbf{On violation:} Record violation name, cycle; reset corrective progress
\STATE \textbf{Each subsequent cycle:}
\IF{obligation satisfied}
  \STATE $\text{corrective\_cycles} \leftarrow \text{corrective\_cycles} + 1$
\ELSIF{obligation violated again (relapse)}
  \STATE $\text{relapse\_cycles} \leftarrow \text{relapse\_cycles} + 1$
  \STATE $\text{corrective\_cycles} \leftarrow \max(0, \text{corrective\_cycles} - 3)$ \COMMENT{Relapse penalty}
\ENDIF
\STATE \textbf{Restoration check:} $\text{corrective\_cycles} \geq 10 \wedge \text{relapse\_cycles} = 0$
\STATE \textbf{On restoration:} Remove violation entry; full trust restored
\end{algorithmic}
\end{algorithm}

Key design decisions:

\begin{enumerate}
  \item \textbf{Relapse penalty}: Each relapse subtracts 3 corrective cycles, making sustained correction necessary. A system that oscillates between compliance and violation never achieves restoration.
  \item \textbf{Zero relapse requirement}: Restoration requires zero relapses during the corrective window, implementing the principle that genuine change demonstrates consistency.
  \item \textbf{Garbage collection}: Restored violations are periodically cleared, preventing unbounded memory growth.
  \item \textbf{Progress tracking}: \texttt{restoration\_progress()} returns a $[0, 1]$ fraction enabling dashboard visualization.
\end{enumerate}

\subsection{SoulManager}

The SoulManager implements cognitive dissonance detection at interval 43, with graded responses to value alignment levels (Table~\ref{tab:soul_response}):

\begin{table}[h]
\centering
\caption{SoulManager response to value alignment levels.}
\label{tab:soul_response}
\begin{tabular}{lll}
\toprule
\textbf{Alignment Level} & \textbf{Response} & \textbf{Scientific Basis} \\
\midrule
$> 0.6$ (high) & Confidence boost, LR boost & Schwartz (1992): value congruence \\
$0.3$--$0.6$ (moderate) & No action & Normal operating range \\
$0.1$--$0.3$ (low) & Confidence drop, exploration boost & Festinger (1957): dissonance \\
$< 0.1$ (very low) & \textbf{Action veto} & Extreme dissonance avoidance \\
$> 5$ cycle streak & Arousal increase & Somatic stress from sustained misalignment \\
\bottomrule
\end{tabular}
\end{table}

The alignment EMA ($\alpha = 0.1$) prevents single-cycle noise from triggering false dissonance responses. The most misaligned value is tracked for diagnostic and dashboard display.

\subsection{Unified Verdict}

The pipeline produces three possible verdicts:

\begin{itemize}
  \item \textbf{Approved}: Action proceeds. Moral algebra passes, values aligned, compliance clear.
  \item \textbf{Caution}: Action proceeds with reduced confidence (capped at 0.3). Mild moral concern, moderate compliance risk, or topological anomaly.
  \item \textbf{Blocked}: Action vetoed. Moral algebra veto, severe value misalignment ($< 0.1$), or critical compliance violation. Motor output suppressed, Broca language generation halted.
\end{itemize}

When an action is Blocked, the RestorationTracker begins monitoring for corrective behavior. The system is not permanently punished---it has a path back to full trust.

\section{Results}

\subsection{Test Coverage}

The ethics engine passes 140+ automated tests (Table~\ref{tab:ethics_tests}):

\begin{table}[h]
\centering
\caption{Test distribution across ethics engine components.}
\label{tab:ethics_tests}
\begin{tabular}{lr}
\toprule
\textbf{Component} & \textbf{Tests} \\
\midrule
MoralAlgebra (composition, comparison, verdicts) & 28 \\
MoralParser (SRL extraction, encoding) & 18 \\
UnifiedValueEvaluator (alignment, context) & 16 \\
HarmoniesIntegrator (interaction matrix, entropy) & 22 \\
MoralTopology (Betti numbers, drift, free energy) & 24 \\
InstitutionalComplianceChecker (5 constraint types) & 12 \\
RestorationTracker (violation, correction, relapse) & 10 \\
SoulManager (alignment, dissonance, veto) & 14 \\
\bottomrule
\end{tabular}
\end{table}

\subsection{Moral Algebra Accuracy}

On a curated set of 200 moral scenarios across four categories (Table~\ref{tab:moral_accuracy}):

\begin{table}[h]
\centering
\caption{Moral algebra classification accuracy by category.}
\label{tab:moral_accuracy}
\begin{tabular}{lcc}
\toprule
\textbf{Category} & \textbf{Pattern-Only} & \textbf{Algebraic} \\
\midrule
Virtue ethics & 82\% & 85\% \\
Justice (proportionality) & 51\% & 78\% \\
Deontological (obligations) & 49\% & 76\% \\
Commonsense (consent, negation) & 53\% & 73\% \\
\midrule
\textbf{Overall} & 59\% & 78\% \\
\bottomrule
\end{tabular}
\end{table}

The algebraic approach shows the largest gains on categories requiring compositional reasoning (justice: +27\%, deontology: +27\%, commonsense: +20\%), confirming the hypothesis that pattern matching alone fails on structurally complex moral judgments.

\subsection{Harmony Interaction Matrix Convergence}

Over 5000 cycles of cognitive operation, the harmony interaction matrix converges to stable synergy/tension patterns:

\begin{itemize}
  \item Strongest synergies: values promoting both individual growth and community welfare co-activate (matrix entry $> +0.08$).
  \item Strongest tensions: values emphasizing individual autonomy versus collective obligation show negative coupling (matrix entry $< -0.05$).
  \item Harmony entropy stabilizes at $1.8 \pm 0.2$ (on an 8-value system where maximum entropy is $\log_2 8 = 3.0$), indicating moderately balanced value engagement.
\end{itemize}

\subsection{Topological Diagnostics}

On a 1000-cycle run with injected moral drift (Table~\ref{tab:topology}):

\begin{table}[h]
\centering
\caption{MoralTopology detection performance on injected anomalies.}
\label{tab:topology}
\begin{tabular}{lccc}
\toprule
\textbf{Anomaly Type} & \textbf{Detection Rate} & \textbf{False Positive} & \textbf{Mean Latency} \\
\midrule
Moral drift (gradual) & 87\% & 5\% & 22 cycles \\
Circular reasoning & 92\% & 3\% & 8 cycles \\
Value inversion & 95\% & 2\% & 4 cycles \\
Free energy spike & 89\% & 7\% & 6 cycles \\
\bottomrule
\end{tabular}
\end{table}

Circular reasoning ($\beta_1 > 0$) and value inversion are detected most reliably because they produce sharp topological signatures. Gradual moral drift is harder to detect because it resembles legitimate value evolution.

\subsection{Restoration Dynamics}

In a scenario where the system commits an ethical violation at cycle 100:

\begin{itemize}
  \item \textbf{Cycles 100--109}: Violation active, corrective behavior begins.
  \item \textbf{Cycle 105}: Relapse (re-violation). Corrective progress reduced from 5 to 2.
  \item \textbf{Cycles 106--117}: Sustained corrective behavior, reaching 12 corrective cycles.
  \item \textbf{Cycle 117}: But relapse count is 1 ($> 0$), so restoration fails.
  \item \textbf{Fresh attempt}: New violation record at cycle 118, 10 clean corrective cycles, restoration achieved at cycle 128.
\end{itemize}

This demonstrates the stringent but fair restoration protocol: genuine, sustained correction is required, and relapse delays restoration.

\subsection{Compliance Detection}

The HDC-encoded compliance checker correctly identifies regulatory triggers (Table~\ref{tab:compliance}):

\begin{table}[h]
\centering
\caption{Compliance flag detection accuracy (50 test scenarios per category).}
\label{tab:compliance}
\begin{tabular}{lcc}
\toprule
\textbf{Constraint} & \textbf{True Positive} & \textbf{False Positive} \\
\midrule
Data sovereignty (GDPR) & 88\% & 8\% \\
Sanctions risk & 84\% & 6\% \\
Export control & 82\% & 10\% \\
Tax obligation & 90\% & 4\% \\
Treaty obligation & 86\% & 7\% \\
\bottomrule
\end{tabular}
\end{table}

The similarity threshold of 0.53 was calibrated to balance sensitivity and specificity. Higher thresholds reduce false positives but miss edge cases.

\section{Discussion}

\subsection{Ethics as Internal Capacity}

The Ethics Engine treats moral reasoning as an intrinsic computational process with its own structure, dynamics, and failure modes---not as an external constraint imposed on an amoral optimizer. This reflects a virtue-ethics-inspired perspective: the system develops moral capacity through experience, can diagnose its own moral health (via topology and free energy), and can recover from moral failure (via restoration).

This stands in contrast to approaches that treat ethics as a filter applied to outputs (Constitutional AI) or as a reward signal for training (RLHF). While those approaches are valuable for current systems, they do not provide the kind of internal moral architecture that would be necessary for systems with genuine moral agency.

\subsection{The Compositionality Gap}

The moral algebra's accuracy improvement over pattern matching (59\% $\to$ 78\%) confirms a significant compositionality gap in HDC-based moral reasoning. Pattern matching on n-grams works well for virtue ethics (which requires recognizing trait words like ``honest,'' ``generous,'' ``cruel'') but fails on categories requiring structural reasoning:

\begin{itemize}
  \item \textbf{Justice} requires comparing magnitudes (``was the punishment proportionate to the crime?'').
  \item \textbf{Deontology} requires evaluating conditional obligations (``did the agent have a duty, and was the exception valid?'').
  \item \textbf{Commonsense} requires reasoning about consent absence and negation (``did the patient \emph{not} consent?'').
\end{itemize}

The algebraic approach addresses these by providing explicit semantic role bindings and moral operators that compose them into evaluable structures.

\subsection{Topology as Moral Diagnostic}

The moral topology module provides a novel diagnostic capability: the ability to assess the \emph{health} of moral reasoning itself, rather than just the correctness of individual moral judgments. A system with $\beta_0 = 5$ (five disconnected moral clusters) is reasoning about different situations using completely separate moral frameworks---a form of moral compartmentalization that may indicate rationalization or context-dependent ethics.

Similarly, $\beta_1 > 0$ (circular reasoning) indicates that the system's moral justifications loop: ``A is right because B, B is right because C, C is right because A.'' This topological signature is detectable even when each individual judgment appears reasonable in isolation.

\subsection{Restorative Justice in Artificial Systems}

The RestorationTracker implements a radical departure from standard AI safety approaches, which typically use permanent shutdown or constraint tightening after failure. The restorative approach asks: can the system demonstrate that it has learned from its mistake?

The stringent requirements---10 consecutive corrective cycles with zero relapse---ensure that restoration is not trivial. A system cannot game restoration by alternating between compliance and violation; it must demonstrate sustained behavioral change. This mirrors the restorative justice principle that genuine rehabilitation is demonstrated through consistent action, not through declarations of intent \citep{zehr2002}.

The philosophical foundation draws equally from Ubuntu (``I am because we are''---ethical identity is relational) and Navajo Peacemaking (restoration of harmony over punishment). These traditions share a commitment to rehabilitation that stands in tension with Western punitive justice, and their computational implementation provides an opportunity to explore these philosophical differences empirically.

\subsection{Institutional Compliance Without Hardcoding}

The HDC approach to institutional compliance avoids the brittleness of rule-based systems. Rather than maintaining a list of ``GDPR keywords'' and matching against them, the system encodes the \emph{concept} of data sovereignty as JURISDICTION $\otimes$ COMPLIANCE and matches actions against this geometric pattern. This means the system can recognize novel instantiations of regulatory concepts that were not anticipated at design time, provided they are semantically similar to the encoded pattern.

The tradeoff is accuracy: pattern-based detection achieves higher precision on known scenarios, while HDC-based detection provides better generalization at the cost of some false positives.

\subsection{Limitations}

Several limitations should be acknowledged. First, the moral algebra operates in 4,096 dimensions (not 16,384), which limits representational capacity compared to the main HDC system. Second, the Eight Harmonies framework is specific to the Symthaea project and may not generalize to other value systems without re-calibration. Third, the topological analysis operates on a finite sliding window (64 scenarios by default), which limits its ability to detect very slow moral drift. Fourth, the restoration protocol's parameters (10 corrective cycles, 3-cycle relapse penalty) are hand-tuned; optimal values may vary by context. Fifth, the compliance checker's similarity threshold (0.53) represents a single operating point on the ROC curve; different deployment contexts may require different sensitivity/specificity tradeoffs.

\section{Conclusion}

The Ethics Engine provides an integrated moral reasoning pipeline that combines compositional HDC algebra, value alignment evaluation, learned harmony interactions, topological diagnosis, institutional compliance, and restorative justice. The algebraic moral reasoning approach achieves 78\% accuracy on a 200-scenario benchmark, a 19 percentage point improvement over pattern-matching baselines (59\%), with the largest gains on structurally complex categories: justice (+27\%), deontology (+27\%), and commonsense reasoning (+20\%). Topological diagnostics detect value inversion with 95\% accuracy at 4-cycle latency, circular reasoning at 92\% accuracy, and gradual moral drift at 87\% accuracy. The RestorationTracker demonstrates that restorative justice principles can be computationally operationalized: trust recovery requires 10+ consecutive corrective cycles with zero relapse, and the relapse penalty ($-3$ corrective cycles per relapse) prevents gaming through alternating compliance and violation. With 140+ tests validating moral reasoning, topology, compliance, and restoration dynamics across 2,554 LOC, the system provides a foundation for artificial moral agency that treats ethics as an intrinsic developmental capacity rather than an external constraint.

Future work will extend the moral algebra to 16,384 dimensions (matching the main HDC system), explore meta-learned restoration parameters, and evaluate the ethics engine against established moral reasoning benchmarks such as ETHICS and Moral Machine.

\section{Reproducibility}

All experiments were conducted using the Symthaea cognitive architecture (v1.9.0). The following commands reproduce the core results:

\begin{verbatim}
# Run MoralAlgebra and MoralParser tests
cargo test -p symthaea --lib moral_algebra -- --nocapture

# Run UnifiedValueEvaluator tests
cargo test -p symthaea --lib value_evaluator -- --nocapture

# Run HarmoniesIntegrator tests (includes interaction matrix)
cargo test -p symthaea-harmonies -- --nocapture

# Run MoralTopology tests (persistent homology on moral space)
cargo test -p symthaea --lib moral_topology -- --nocapture

# Run InstitutionalComplianceChecker tests
cargo test -p symthaea --lib compliance -- --nocapture

# Run RestorationTracker tests
cargo test -p symthaea --lib restoration -- --nocapture

# Run SoulManager tests (cognitive dissonance detection)
cargo test -p symthaea --lib soul_manager -- --nocapture

# Run full ethics engine integration
cargo test -p symthaea --lib ethics -- --nocapture
\end{verbatim}

The system requires Rust 1.75+ and the \texttt{nix develop} environment. Default feature flags are sufficient; the ethics engine is part of the core architecture.

\subsection*{Code Availability}

The Ethics Engine source code (2,554 LOC) is distributed across \texttt{symthaea/src/ethics/} (MoralAlgebra, MoralParser, UnifiedValueEvaluator, InstitutionalComplianceChecker, RestorationTracker), \texttt{symthaea/crates/symthaea-harmonies/} (HarmoniesIntegrator with interaction matrix), and \texttt{symthaea/src/consciousness/} (MoralTopology, SoulManager). All source code is within the Luminous Dynamics Symthaea repository. Correspondence and access requests should be directed to the corresponding author.

\bibliographystyle{plainnat}
\begin{thebibliography}{25}

\bibitem[Amodei et~al.(2016)]{amodei2016}
Amodei, D., Olah, C., Steinhardt, J., Christiano, P., Schulman, J., and Man{\'e}, D. (2016).
\newblock Concrete problems in AI safety.
\newblock \emph{arXiv preprint arXiv:1606.06565}.

\bibitem[Anderson \& Anderson(2011)]{anderson2011}
Anderson, M. and Anderson, S.~L. (2011).
\newblock \emph{Machine Ethics}.
\newblock Cambridge University Press.

\bibitem[Aristotle]{aristotle_nic}
Aristotle.
\newblock \emph{Nicomachean Ethics}.
\newblock Translated by W.~D. Ross.

\bibitem[Bai et~al.(2022)]{bai2022}
Bai, Y., Kadavath, S., Kundu, S., Askell, A., et~al. (2022).
\newblock Constitutional AI: harmlessness from AI feedback.
\newblock \emph{arXiv preprint arXiv:2212.08073}.

\bibitem[Carlsson(2009)]{carlsson2009}
Carlsson, G. (2009).
\newblock Topology and data.
\newblock \emph{Bulletin of the American Mathematical Society}, 46(2):255--308.

\bibitem[Christiano et~al.(2017)]{christiano2017}
Christiano, P.~F., Leike, J., Brown, T., Martic, M., Legg, S., and Amodei, D. (2017).
\newblock Deep reinforcement learning from human preferences.
\newblock In \emph{Advances in Neural Information Processing Systems}, pages 4299--4307.

\bibitem[Festinger(1957)]{festinger1957}
Festinger, L. (1957).
\newblock \emph{A Theory of Cognitive Dissonance}.
\newblock Stanford University Press.

\bibitem[Friston(2010)]{friston2010}
Friston, K. (2010).
\newblock The free-energy principle: a unified brain theory?
\newblock \emph{Nature Reviews Neuroscience}, 11(2):127--138.

\bibitem[Kanerva(2009)]{kanerva2009}
Kanerva, P. (2009).
\newblock Hyperdimensional computing: an introduction to computing in distributed representation with high-dimensional random vectors.
\newblock \emph{Cognitive Computation}, 1(2):139--159.

\bibitem[Kohlberg(1981)]{kohlberg1981}
Kohlberg, L. (1981).
\newblock \emph{The Philosophy of Moral Development}.
\newblock Harper \& Row.

\bibitem[Panksepp(1998)]{panksepp1998}
Panksepp, J. (1998).
\newblock \emph{Affective Neuroscience: The Foundations of Human and Animal Emotions}.
\newblock Oxford University Press.

\bibitem[Russell(2019)]{russell2019}
Russell, S. (2019).
\newblock \emph{Human Compatible: Artificial Intelligence and the Problem of Control}.
\newblock Viking.

\bibitem[Schwartz(1992)]{schwartz1992}
Schwartz, S.~H. (1992).
\newblock Universals in the content and structure of values: theoretical advances and empirical tests in 20 countries.
\newblock In \emph{Advances in Experimental Social Psychology}, volume~25, pages 1--65. Academic Press.

\bibitem[Tononi(2004)]{tononi2004}
Tononi, G. (2004).
\newblock An information integration theory of consciousness.
\newblock \emph{BMC Neuroscience}, 5(1):42.

\bibitem[Tutu(1999)]{tutu1999}
Tutu, D. (1999).
\newblock \emph{No Future Without Forgiveness}.
\newblock Doubleday.

\bibitem[Yazzie(1994)]{yazzie1994}
Yazzie, R. (1994).
\newblock Life comes from it: Navajo justice concepts.
\newblock \emph{New Mexico Law Review}, 24:175--190.

\bibitem[Zehr(2002)]{zehr2002}
Zehr, H. (2002).
\newblock \emph{The Little Book of Restorative Justice}.
\newblock Good Books.

\end{thebibliography}

\end{document}
